%==============================================================================
% APPENDICES
%==============================================================================
\appendix
\chapter*{Appendices}
\addcontentsline{toc}{chapter}{Appendices}
\markboth{Appendices}{Appendices}

\renewcommand{\thesection}{Appendix \Alph{section}}
\setcounter{section}{0}
\addtocontents{toc}{\protect\tocAppendixWidth}

%---------------------------------------------------------------- Appendix A --
\section{Software environment and installation procedure}
\label{ann:install}

\noindent
\begin{tabularx}{\textwidth}{@{}lX@{}}
\toprule
\textbf{Component} & \textbf{Version / role}\\
\midrule
Operating system & Linux (Ubuntu 24.04) under WSL2\\
OpenFOAM & v2506 (OpenCFD distribution)\\
OFFBEAT & \textit{master} branch, compiled from source\\
Python & 3.12, dedicated virtual environment\\
Language model & \num{7} billion parameters, run locally\\
Embedding model & local \textbf{multilingual} model (see §\ref{sec:rag})\\
Post-processing & VTK reading library\\
Surrogate & Gaussian process (Matérn kernel)\\
\bottomrule
\end{tabularx}

\vspace{0.5cm}

\noindent
\textit{Remark.} The project and the simulation cases must reside on the native
Linux file system: the input/output of an OpenFOAM case is severely degraded on
a file system mounted from the host.

%---------------------------------------------------------------- Appendix B --
\section{Structure of an OFFBEAT case and exposed parameters}
\label{ann:cas}

\noindent An OFFBEAT case is a file tree:

\vspace{0.3cm}
\noindent
\begin{tabularx}{\textwidth}{@{}lX@{}}
\toprule
\textbf{Element} & \textbf{Contents}\\
\midrule
\texttt{0/} & Initial fields (temperature, displacement, neutron flux)\\
\texttt{constant/} & Material properties, physical models, power history,
  coolant pressure\\
\texttt{system/} & Run control (duration, time step, writing), numerical
  schemes and solvers\\
\texttt{rodDict} & Rod geometry by blocks (radii, heights, discretisation),
  specific to OFFBEAT\\
\bottomrule
\end{tabularx}

\vspace{0.6cm}

\noindent Parameters exposed by the agent:

\vspace{0.3cm}
\noindent
\begin{tabularx}{\textwidth}{@{}lXl@{}}
\toprule
\textbf{Parameter} & \textbf{Meaning} & \textbf{Unit}\\
\midrule
\texttt{linear\_heat\_rate} & Linear heat rate & \si{\watt\per\metre}\\
\texttt{end\_time} & Simulated duration & \si{\second}\\
\texttt{delta\_t} & Initial time step & \si{\second}\\
\texttt{write\_interval} & Writing frequency & steps\\
\texttt{enrichment} & Enrichment & --\\
\texttt{fuel\_outer\_radius} & Pellet outer radius & \si{\milli\metre}\\
\texttt{fuel\_height} & Height of the column & \si{\milli\metre}\\
\texttt{clad\_inner\_radius} & Cladding inner radius & \si{\milli\metre}\\
\texttt{clad\_outer\_radius} & Cladding outer radius & \si{\milli\metre}\\
\texttt{n\_cells\_r\_fuel} & Radial discretisation & --\\
\bottomrule
\end{tabularx}

\vspace{0.4cm}
\noindent
\textit{Consistency constraint checked before execution:}
$r_{\text{pellet}} < r_{\text{clad,in}} < r_{\text{clad,out}}$.

%---------------------------------------------------------------- Appendix C --
\section{Error knowledge base}
\label{ann:errkb}

\noindent Each entry associates a detection pattern with a diagnosis and a
correction. Extract of the empirically validated entries:

\vspace{0.3cm}
\noindent
\begin{tabularx}{\textwidth}{@{}>{\raggedright\arraybackslash}p{3.4cm}XX@{}}
\toprule
\textbf{Identifier} & \textbf{Diagnosis} & \textbf{Correction}\\
\midrule
Floating-point exception & Division by zero or non-numeric value: divergence,
  unrealistic conditions, or a material property out of range
  & Reduce the time step\\
Duration beyond the power history & The simulated duration exceeds the last
  tabulated point of the history & Read the time tables present in the case and
  bring the duration just below the last point\\
Restart impossible & Boundary condition not re-readable from a written time
  step & Clean the time steps and restart cleanly\\
Inverted geometry & Invalid mesh & None: intervention required\\
Missing dictionary & Entry absent from a configuration file &
  None: intervention required\\
\bottomrule
\end{tabularx}

\vspace{0.4cm}
\noindent\textit{Remark.} An entry inherited from the reference project,
relating to the Courant number, turned out to be irrelevant to OFFBEAT: the
solver contains no advection term, hence no stability condition of that kind. It
was kept disabled, together with its justification, as a record.

%---------------------------------------------------------------- Appendix D --
\section{Safety criteria knowledge base}
\label{ann:safetykb}

\noindent
\begin{tabularx}{\textwidth}{@{}>{\raggedright\arraybackslash}p{3.3cm}
                                >{\raggedright\arraybackslash}p{2.2cm}
                                >{\raggedright\arraybackslash}p{1.9cm}
                                X@{}}
\toprule
\textbf{Criterion} & \textbf{Quantity} & \textbf{Limit} & \textbf{Remark}\\
\midrule
Centreline melting & Temperature, maximum & \SI{3113}{\kelvin} &
  Unirradiated $\mathrm{UO_2}$; the limit decreases with burnup (see
  below)\\
PCMI strain & \textbf{Plastic} $\theta\theta$ strain, cladding zone &
  \SI{1}{\percent} & Usual design criterion; see [7], [8]\\
Permanent strain & Equivalent creep strain, cladding zone &
  \SI{1}{\percent} & Limit taken over by default, \textbf{not validated}\\
Cladding stress & $\theta\theta$ stress, cladding zone &
  \texttt{sigmaY} field & Limit read from the field computed by the solver
  (fallback \SI{250}{\mega\pascal})\\
Gap closure & Gap width, minimum & \SI{5}{\micro\metre} &
  Reversed sense: the danger is a value that is too \emph{small}\\
\bottomrule
\end{tabularx}

\vspace{0.4cm}
\noindent
\textcolor{cAlert}{\textbf{Important.}} These five criteria are marked as
\textbf{not validated} in the base. They are published design orders of
magnitude, not invented values, but their validation with the supervisor remains
a prerequisite to any exploitation.

\vspace{0.4cm}
\noindent
\textbf{The decrease of the melting limit with burnup.} The melting temperature
of $\mathrm{UO_2}$ decreases during irradiation: the fission products and the
plutonium formed enter into solid solution in the matrix and lower the melting
point, as a solute lowers that of a solvent. The margin to melting therefore
narrows over the cycle, at the very moment when the centreline temperature is
itself rising.

The note initially carried in the base indicated a slope of \SI{-32}{\kelvin}
per \num{10}~GWd/tU. The bibliographic check carried out for this report showed
that there are in fact \textbf{two published correlations}, and that the one
adopted by the project was the older of the two (table~\ref{tab:fusion}).

\begin{table}[htbp]
\centering
\caption[UO2 melting correlations]{The two published correlations for the
decrease of the melting temperature of $\mathrm{UO_2}$ with burnup. Both start
from \SI{3113.15}{\kelvin} for unirradiated $\mathrm{UO_2}$ (Brassfield, 1968).}
\label{tab:fusion}
\small
\begin{tabularx}{\textwidth}{@{}Xccc@{}}
\toprule
\textbf{Correlation} & \textbf{Subroutine} & \textbf{Slope} &
\textbf{$T_{\text{melt}}$ at 50 GWd/tU}\\
\midrule
MATPRO (earlier) & \texttt{FHYPRP} & $-3.2$ K/GWd/tU
  & \SI{2953}{\kelvin}\\
FRAPCON-4.0 / FRAPTRAN-2.0 (adopted) & \texttt{PHYPRP} &
  $-0.5$ K/GWd/tU & \SI{3088}{\kelvin}\\
\bottomrule
\end{tabularx}
\end{table}
